%%%%%%%%%%%%%%%%%%%%%%%%%%%%%%%%%%%%%%%%%
% Masters/Doctoral Thesis 
% LaTeX Template
% Version 2.4 (22/11/16)
%
% This template has been downloaded from:
% http://www.LaTeXTemplates.com
%
% Version 2.x major modifications by:
% Vel (vel@latextemplates.com)
%
% This template is based on a template by:
% Steve Gunn (http://users.ecs.soton.ac.uk/srg/softwaretools/document/templates/)
% Sunil Patel (http://www.sunilpatel.co.uk/thesis-template/)
%
% Template license:
% CC BY-NC-SA 3.0 (http://creativecommons.org/licenses/by-nc-sa/3.0/)
%
%%%%%%%%%%%%%%%%%%%%%%%%%%%%%%%%%%%%%%%%%

%----------------------------------------------------------------------------------------
%	PACKAGES AND OTHER DOCUMENT CONFIGURATIONS
%----------------------------------------------------------------------------------------

\documentclass[
11pt, % The default document font size, options: 10pt, 11pt, 12pt
%oneside, % Two side (alternating margins) for binding by default, uncomment to switch to one side
english, % ngerman for German
singlespacing, % Single line spacing, alternatives: onehalfspacing or doublespacing
%draft, % Uncomment to enable draft mode (no pictures, no links, overfull hboxes indicated)
%nolistspacing, % If the document is onehalfspacing or doublespacing, uncomment this to set spacing in lists to single
%liststotoc, % Uncomment to add the list of figures/tables/etc to the table of contents
%toctotoc, % Uncomment to add the main table of contents to the table of contents
%parskip, % Uncomment to add space between paragraphs
%nohyperref, % Uncomment to not load the hyperref package
headsepline, % Uncomment to get a line under the header
%chapterinoneline, % Uncomment to place the chapter title next to the number on one line
%consistentlayout, % Uncomment to change the layout of the declaration, abstract and acknowledgements pages to match the default layout
]{MastersDoctoralThesis} % The class file specifying the document structure

\usepackage[utf8]{inputenc} % Required for inputting international characters
\usepackage[T1]{fontenc} % Output font encoding for international characters

\usepackage{palatino} % Use the Palatino font by default

\usepackage[backend=bibtex,style=authoryear,natbib=true]{biblatex} % Use the bibtex backend with the authoryear citation style (which resembles APA)

\addbibresource{example.bib} % The filename of the bibliography

\usepackage[autostyle=true]{csquotes} % Required to generate language-dependent quotes in the bibliography

\usepackage{graphicx}

\usepackage{amsmath}
\usepackage{algorithm}
\usepackage[noend]{algpseudocode}

\makeatletter
\def\BState{\State\hskip-\ALG@thistlm}
\makeatother

\floatname{algorithm}{Procedure}
\renewcommand{\algorithmicrequire}{\textbf{Input:}}
\renewcommand{\algorithmicensure}{\textbf{Output:}}
\usepackage{textcomp}

%----------------------------------------------------------------------------------------
%	MARGIN SETTINGS
%----------------------------------------------------------------------------------------

\geometry{
	paper=a4paper, % Change to letterpaper for US letter
	inner=2.5cm, % Inner margin
	outer=3.8cm, % Outer margin
	bindingoffset=.5cm, % Binding offset
	top=1.5cm, % Top margin
	bottom=1.5cm, % Bottom margin
	%showframe, % Uncomment to show how the type block is set on the page
}

%----------------------------------------------------------------------------------------
%	THESIS INFORMATION
%----------------------------------------------------------------------------------------

\thesistitle{Computational Support for Concept Blending applied to Musical Instruments} % Your thesis title, this is used in the title and abstract, print it elsewhere with \ttitle
\supervisor{Prof. Bjørnar \textsc{Tessem}} % Your supervisor's name, this is used in the title page, print it elsewhere with \supname
\examiner{} % Your examiner's name, this is not currently used anywhere in the template, print it elsewhere with \examname
\degree{Master} % Your degree name, this is used in the title page and abstract, print it elsewhere with \degreename
\author{Andreas \textsc{Naustdal}} % Your name, this is used in the title page and abstract, print it elsewhere with \authorname
\addresses{} % Your address, this is not currently used anywhere in the template, print it elsewhere with \addressname

\subject{Information Science} % Your subject area, this is not currently used anywhere in the template, print it elsewhere with \subjectname
\keywords{} % Keywords for your thesis, this is not currently used anywhere in the template, print it elsewhere with \keywordnames
\university{\href{http://www.uib.no}{University Of Bergen}} % Your university's name and URL, this is used in the title page and abstract, print it elsewhere with \univname
\department{\href{http://www.uib.no/infomedia}{Department of Information Science and Media Studies}} % Your department's name and URL, this is used in the title page and abstract, print it elsewhere with \deptname
\group{\href{http://researchgroup.university.com}{Research Group Name}} % Your research group's name and URL, this is used in the title page, print it elsewhere with \groupname
\faculty{\href{http://www.uib.no/svf}{Faculty of Social Sciences}} % Your faculty's name and URL, this is used in the title page and abstract, print it elsewhere with \facname

\AtBeginDocument{
\hypersetup{pdftitle=\ttitle} % Set the PDF's title to your title
\hypersetup{pdfauthor=\authorname} % Set the PDF's author to your name
\hypersetup{pdfkeywords=\keywordnames} % Set the PDF's keywords to your keywords
}

\begin{document}

\frontmatter % Use roman page numbering style (i, ii, iii, iv...) for the pre-content pages

\pagestyle{plain} % Default to the plain heading style until the thesis style is called for the body content

%----------------------------------------------------------------------------------------
%	TITLE PAGE
%----------------------------------------------------------------------------------------

\begin{titlepage}
\begin{center}

\vspace*{.06\textheight}
{\scshape\LARGE \univname\par}\vspace{1.5cm} % University name
\textsc{\Large Masters Thesis}\\[0.5cm] % Thesis type

\HRule \\[0.4cm] % Horizontal line
{\huge \bfseries \ttitle\par}\vspace{0.4cm} % Thesis title
\HRule \\[1.5cm] % Horizontal line
 
\begin{minipage}[t]{0.4\textwidth}
\begin{flushleft} \large
\emph{Author:}\\
{\authorname} % Author name - remove the \href bracket to remove the link
\end{flushleft}
\end{minipage}
\begin{minipage}[t]{0.4\textwidth}
\begin{flushright} \large
\emph{Supervisor:} \\
{\supname} % Supervisor name - remove the \href bracket to remove the link  
\end{flushright}
\end{minipage}\\[3cm]
 
\vfill

\large \textit{A thesis submitted in fulfillment of the requirements\\ for the degree of \degreename}\\[0.3cm] % University requirement text
\textit{in the}\\[0.4cm]
\deptname\\[2cm] % Research group name and department name
 
\vfill

{\large \today}\\[4cm] % Date
%\includegraphics{Logo} % University/department logo - uncomment to place it

{\large Version 0.3}\\[4cm] % Date
 
\vfill
\end{center}
\end{titlepage}

%----------------------------------------------------------------------------------------
%	DECLARATION PAGE
%----------------------------------------------------------------------------------------

%\begin{declaration}
%\addchaptertocentry{\authorshipname} % Add the declaration to the table of contents
%\noindent I, \authorname, declare that this thesis titled, \enquote{\ttitle} and the work presented in it are my own. I confirm that:

%\begin{itemize} 
%\item This work was done wholly or mainly while in candidature for a research degree at this University.
%\item Where any part of this thesis has previously been submitted for a degree or any other qualification at this University or any other institution, this has been clearly stated.
%\item Where I have consulted the published work of others, this is always clearly attributed.
%\item Where I have quoted from the work of others, the source is always given. With the exception of such quotations, this thesis is entirely my own work.
%\item I have acknowledged all main sources of help.
%\item Where the thesis is based on work done by myself jointly with others, I have made clear exactly what was done by others and what I have contributed myself.\\
%\end{itemize}
 
%\noindent Signed:\\
%\rule[0.5em]{25em}{0.5pt} % This prints a line for the signature
 
%\noindent Date:\\
%\rule[0.5em]{25em}{0.5pt} % This prints a line to write the date
%\end{declaration}

%\cleardoublepage

%----------------------------------------------------------------------------------------
%	QUOTATION PAGE
%----------------------------------------------------------------------------------------

%\vspace*{0.2\textheight}

%\noindent\enquote{\itshape Thanks to my solid academic training, today I can write hundreds of words on virtually any topic without possessing a shred of information, which is how I got a good job in journalism.}\bigbreak

%\hfill Dave Barry

%----------------------------------------------------------------------------------------
%	ABSTRACT PAGE
%----------------------------------------------------------------------------------------

\begin{abstract}
\addchaptertocentry{\abstractname} % Add the abstract to the table of contents
This thesis provides a concept blending implementation that suggests which properties of a known concept are most compatible to blend with another concept. The implementation uses WikiPedia descriptions of concepts as data source, and NLP tools such as WordNet and stanford CoreNLP to prepare a representation of concepts with properties. By joining the generalized properties of two concepts in a tree structure, we look for patterns in the data which correlate with properties that make sense to blend. A heuristic function based on these patterns is used to rank the subtrees to return top suggestions of features to blend between the two concepts.
%What are the characteristics of elements in a concept that are suitable for concept blending?
%Can a concept blending algorithm generate musical instruments that are reasonable and useful?
%What are optimal solutions for preparing input data and perform contradiction detection, in the context of using the concept blending algorithm to generate musical instruments?
%What is the expectations of this thesis?
%Results
%The Thesis Abstract is written here (and usually kept to just this page). The page is kept centered vertically so can expand into the blank space above the title too\ldots
\end{abstract}

%----------------------------------------------------------------------------------------
%	ACKNOWLEDGEMENTS
%----------------------------------------------------------------------------------------

\begin{acknowledgements}
\addchaptertocentry{\acknowledgementname} % Add the acknowledgements to the table of contents
I would like to thank Prof. Bjørnar Tessem for supervising and supporting the thesis. I would like to thank Dr. Ankica Babic for helpful feedback and guidance in creating a project outline
\ldots
\end{acknowledgements}

%----------------------------------------------------------------------------------------
%	LIST OF CONTENTS/FIGURES/TABLES PAGES
%----------------------------------------------------------------------------------------

\tableofcontents % Prints the main table of contents

\listoffigures % Prints the list of figures

\listoftables % Prints the list of tables

%----------------------------------------------------------------------------------------
%	ABBREVIATIONS
%----------------------------------------------------------------------------------------

\begin{abbreviations}{ll} % Include a list of abbreviations (a table of two columns)

\textbf{POS} & \textbf{P}art \textbf{O}f \textbf{S}peech\\
\textbf{D} & \textbf{D}epth of word in tree\\
\textbf{BL} & \textbf{B}ranch \textbf{L}ength from shared hypernym\\

\end{abbreviations}

%----------------------------------------------------------------------------------------
%	PHYSICAL CONSTANTS/OTHER DEFINITIONS
%----------------------------------------------------------------------------------------

%\begin{constants}{lr@{${}={}$}l} % The list of physical constants is a three column table

% The \SI{}{} command is provided by the siunitx package, see its documentation for instructions on how to use it

%Speed of Light & $c_{0}$ & \SI{2.99792458e8}{\meter\per\second} (exact)\\
%Constant Name & $Symbol$ & $Constant Value$ with units\\

%\end{constants}

%----------------------------------------------------------------------------------------
%	SYMBOLS
%----------------------------------------------------------------------------------------

%\begin{symbols}{lll} % Include a list of Symbols (a three column table)

%$a$ & distance & \si{\meter} \\
%$P$ & power & \si{\watt} (\si{\joule\per\second}) \\
%Symbol & Name & Unit \\

%\addlinespace % Gap to separate the Roman symbols from the Greek

%$\omega$ & angular frequency & \si{\radian} \\

%\end{symbols}

%----------------------------------------------------------------------------------------
%	DEDICATION
%----------------------------------------------------------------------------------------

%\dedicatory{For/Dedicated to/To my\ldots} 

%----------------------------------------------------------------------------------------
%	THESIS CONTENT - CHAPTERS
%----------------------------------------------------------------------------------------

\mainmatter % Begin numeric (1,2,3...) page numbering

\pagestyle{thesis} % Return the page headers back to the "thesis" style

% Include the chapters of the thesis as separate files from the Chapters folder
% Uncomment the lines as you write the chapters

% Chapter 1

\chapter{Introduction} % Main chapter title

\label{Chapter1} % For referencing the chapter elsewhere, use \ref{Chapter1} 

%----------------------------------------------------------------------------------------

% Define some commands to keep the formatting separated from the content 
\newcommand{\keyword}[1]{\textbf{#1}}
\newcommand{\tabhead}[1]{\textbf{#1}}
\newcommand{\code}[1]{\texttt{#1}}
\newcommand{\file}[1]{\texttt{\bfseries#1}}
\newcommand{\option}[1]{\texttt{\itshape#1}}

%----------------------------------------------------------------------------------------

\section{Motivation/Background}
Combining two concepts into a new concept that combines elements from each original concept is something humans can do rather easily. If we think of blending a horse and a bird, we can easily imagine a winged horse like the mythological Pegasus. Or when blending a car with a boat we can think of an amphibian boat. Formalizing the task of concept blending is not as straightforward. The problem consists of retrieving input data, selecting which parts of the data that should be blended and handling contradictory elements in the final blend.

For the concept blending algorithm to be useful, it is ideal that the pieces that are merging between the concepts make sense to join. In order to do this, we want to find a way to decompose the concepts into properties containing enough information so that we can look for patterns to identify suitable blending suggestions. %This step of concept blending is what we are going to focus on in this thesis.
\\The goal of this master thesis is to implement an algorithm to provide suggestions of elements from one concept that are suitable for transfer to another. Descriptions of the two concepts are retrieved from a public data source, DBpedia, and natural language processing tools like WordNet and Stanford Parser are used to find which elements the concept consists of. Although the concept blending algorithm may be used on any concepts, I will focus mainly on a certain domain such as musical instruments, with the goal of designing new instruments. Further the generated instruments will be evaluated on how useful they are in a musical context. The goal is to find out whether this approach can produce new instruments that make sense and find out how the algorithm can be useful to musicians and instrument makers.

%----------------------------------------------------------------------------------------

\section{Research Questions}
Here we list three questions that we want to research through incremental improvements and evaluation of the algorithm.
\begin{itemize}
\item What are the characteristics of elements in a concept that are suitable for concept blending?
\end{itemize}
The reason we want to find these are because the characteristics are what we can use as a base for the parameters in the heuristic algorithm so we can rank the good candidates for blending in any concept.
\begin{itemize}
\item How should we represent the concepts in order to find patterns of good elements for blending and detect contradictions indicating bad blends?
\end{itemize}
We want to find out which data are optimal for the input stage in order for the algorithm to work optimally and give the best suggestions. Can the concepts and their parts keep their original meaning and relations in the final results?
\begin{itemize}
\item Can a concept blending algorithm generate musical instruments that are reasonable and useful?
\end{itemize}
We want to evaluate the output of the algorithm to see if it is valuable in its environment.
%\item What is the optimal solution for preparing input data and detect contradictions, when using the concept blending algorithm to generate musical instruments?

The motivation for these questions is to uncover new knowledge on what the concept blending algorithm can be used for, how it can be used successfully and how it compares to other approaches.

%----------------------------------------------------------------------------------------
% Chapter 2

\chapter{Literature review} % Main chapter title

\label{Chapter2} % For referencing the chapter elsewhere, use \ref{Chapter2} 

%----------------------------------------------------------------------------------------

\section{Creativity}
%When looking at different definitions of creativity, it is quite abstract. Here are some definitions of creativity:

%Wikipedia:
%Creativity is a phenomenon whereby something new and somehow valuable is formed. The created item may be intangible (such as an idea, a scientific theory, a musical composition, or a joke) or a physical object (such as an invention, a literary work, or a painting).

%Google/Oxford:
%the use of imagination or original ideas to create something; inventiveness.
%"firms are keen to encourage creativity"
%synonyms:
%inventiveness, imagination, imaginativeness, innovation, innovativeness, originality, individuality; 
%artistry, expressiveness, inspiration, vision, creative power, creative talent, creative gift, creative skill, resourcefulness, ingenuity, enterprise
%"challenging objectives motivate staff and encourage creativity"


Cambridge Dictionary gives the following definition of creativity:
the ability to produce original and unusual ideas, or to make something new or imaginative.

%Dictionary.com
%noun
%1.
%the state or quality of being creative.
%2.
%the ability to transcend traditional ideas, rules, patterns, relationships, or the like, and to create meaningful new ideas, forms, methods, interpretations, etc.; originality, progressiveness, or imagination:
%the need for creativity in modern industry; creativity in the performing arts.
%3.
%the process by which one utilizes creative ability:
%Extensive reading stimulated his creativity.

%Merriam-Webster:
%Definition of creativity
%1:  the quality of being creative
%2:  the ability to create her artistic creativity

Creativity happens in different ways. Sometimes a new invention is found because the right elements were available, and a combination of these created something new. For instance the first electric light was invented when Humphry Davy had access to the world's largest battery at the time, and passing current through a strip of platinum produced light. In other cases creativity is goal-driven. When electric light was commercialized it went through a creative process of choosing the right materials and finding an cost-efficient production method with the goal of making the light-bulb profitable.
%It seems hard to define specifically how creativity is actually done, since it can be done in many ways. 
%The established literature seems to branch into different methods. Examples of creative tasks include invention, problem solving and art. Our approach will follow the theory of conceptual blending.

\parencite{busse1980theories} lists seven categories of creativity theories: psychoanalytic, gestalt, association, perceptual, humanistic, cognitive-developmental and composite theories.
They also draw a line between convergent and divergent problems that the theories are applicable to. The difference is convergent problems have one or few right answers while divergent has many possible solutions. One end is a more goal-driven and problem-solving type of creativity and the other end includes more open-ended activities such as art and conceiving novel ideas.

In the association theories, creativity is thought to come from associations. Arthur Koestler combines associationism and psychoanalytic concepts in what he called the \emph{bisociation} theory \parencite{busse1980theories}. In bisociation, two independent matrices of ideas are combined, guided by subconscious processes. Bisociation inspired the theory of conceptual blending by Gilles Fauconnier and Mark Turner \parencite{fauconnier2002way}.

%----------------------------------------------------------------------------------------

%\section{Computational creativity}


%----------------------------------------------------------------------------------------

\section{Conceptual blending}
Conceptual blending is a theory proposed by Gilles Fauconnier and Mark Turner. It works by projecting structure from input mental spaces to a separate, "blended" mental space. \parencite{reference3}
%According to \parencite{reference3}: “concept blending constitutes a cognitive process which allows for the combination of certain elements (and their relations) from originally distinct conceptual spaces into a new unified space combining these previously separate elements and allowing the performance of reasoning and inference over the combination.”
\begin{figure}
\centering
\includegraphics[width=0.7\linewidth]{"../../../Desktop/Master/concept blending firkant"}
\caption{Conceptual blending model showing mental spaces (circles) and mappings of counterpart connections (solid lines)}
\label{fig:concept-blending-firkant}
\end{figure}


%----------------------------------------------------------------------------------------

\section{Analogy}

%Relasjonar er viktig i analogiar
%Analogiar er viktig for concept blending
%Generalisering?

%Dedre Gentner 1983

%Structure mapping engine is a direct mapping approach

%Generalization based models of analogy
%Smaling 2003
%%HDTP (Schmidt et al. 2014)

%Min løysing er delvis inspirert av analogi ved generalisering?

%----------------------------------------------------------------------------------------

\section{Amalgam}
One approach to concept blending is the use of amalgams. \parencite{Reference4}
An amalgam uses generalization of two concepts to make them compatible for blending through unification of the two generalizations. When parts of the concepts share a similar generalization, but are different, they are not unifiable. But they can unify with the \emph{least general generalization}. To make the concepts unifiable, we can therefore generalize corresponding pairs of parts between the concepts. For instance, when creating an amalgam of a \emph{red French vehicle} and a \emph{German minivan}, we can make them unifiable by generalizing the first to a \emph{red European vehicle}. When unifying we get a \emph{red German minivan}. Or if we generalize the second to a \emph{European minivan}, we can unify to a \emph{red French minivan}.
%An amalgam of two descriptions is a new description that contains parts from these two descriptions. (sitat frå Besold eller Ontanón and Plaza 2010?)

%Subsumption relasjon (hyponym subsuming hypernym)
%All informasjon i hypernymet er sant for hyponymet

%----------------------------------------------------------------------------------------

\section{Meaning and context}
An important part of concept blending is not only representing the parts to blend as words, but also their meaning which depends on the context. \parencite{Reference6} propose that there is a difference between lexical concepts and meaning. Meaning is not a function of language, but arise from the use of language. \parencite{Reference6} provides a theory of lexical concept integration, the Theory of Lexical Concepts and Cognitive Models (LCCM). 

%----------------------------------------------------------------------------------------

%\section{The COINVENT Project}

%---------------------------- 
% Chapter 3

\chapter{Methods and methodologies} % Main chapter title

\label{Chapter3} % For referencing the chapter elsewhere, use \ref{Chapter3} 

%----------------------------------------------------------------------------------------

\section{Personal Kanban}
In order to keep my work organized and productive, I used the agile method called Personal Kanban. To visualize the workload, Personal Kanban uses a value stream which are columns that divide the work based on the progress of each task. I used the simplest one that divides the work into three columns. The first is the project backlog which includes all the tasks that awaits production. The middle column is the tasks currently being worked on, which should be a limited number of tasks so that they are not getting stuck half-finished and distracting the task one is actually doing. The last column is the tasks that have been completely done so that there are no reason to be distracted by it anymore. \parencite{Reference1}
\\\\
To keep a workflow where I can spend my energy on making sure the model, view and controller components is interacting well together, I limited my task number to 3. This way I could select one task from each component that depended on each other, while unfinished tasks would not be slowing down the workflow.\\
To review my progress, I had a daily evaluation inspired by the stand-up meeting used in XP \parencite{Reference2} where I answered the following questions:
\begin{enumerate}
\item What was accomplished yesterday?
\item What will be attempted today?
\item What problems are causing delays?
\end{enumerate}


%----------------------------------------------------------------------------------------

\section{Proof of concept}

%----------------------------------------------------------------------------------------
% Chapter 4

\chapter{Development} % Main chapter title

\label{Chapter4} % For referencing the chapter elsewhere, use \ref{Chapter4} 

%----------------------------------------------------------------------------------------

\section{Tools and technologies}
In this thesis we have developed an algorithm which provide suggestions of which properties that are most compatible for blending between two concepts of a domain such as musical instruments. To achieve this we have used natural language processing tools such as Stanford coreNLP and WordNet.

\subsection{WikiPedia}
WikiPedia is a large online encyclopedia, consisting of over 5 million english articles. We use it as a data source for the concept descriptions. We query WikiPedia for the concept or instrument we want to blend. Then we take the first section in the article, which we refer to as the abstract. The results depend on the quality of this abstract. Sometimes it includes irrelevant words, for example cultural references such as the inventor or where the instrument is produced.

\subsection{WordNet}
WordNet contains a large group of English words that are linked together with its synonym group called synset. Each synset has a distinct meaning, so a word that can mean different things has one synset for each meaning. This is useful for preparing the data so we do not include two elements that mean the same. WordNet also provides antonyms of a word, which is useful to detect contradictions in the blended concept.

\subsection{Stanford coreNLP toolkit}
The Stanford coreNLP toolkit is widely used for natural language analysis. It includes tools for tokenization, sentence splitting, part-of-speech tagging, morphological analysis, named entity recognition, syntactic tagging, coreference resolution and more. \parencite{manning-EtAl:2014:P14-5}

\subsection{JavaScript Frameworks and Libraries}

The implementation was written in JavaScript using the frameworks AngularJS and Node.JS. The algorithm uses the following libraries from the Node Package Manager (NPM):
\subsection{wordnet-magic}
wordnet-magic is a Node.JS implementation of Princeton's WordNet lexical database for the English language. We use it to retrieve a set of synsets for each word in the description of the concepts. When the correct synsets are chosen, we also use it to retrieve all hypernyms of the synset in a tree structure where every ancestor is stored as synsets. 

\subsection{stanford-simple-nlp}
stanford-simple-nlp is a Node.JS wrapper for StanfordCoreNLP, a set of natural language processing tools. We use two of their tools, a sentence parser and a part-of-speech tagger. We use the sentence parser to identify noun phrases when we want to find synsets of elements that are represented in multiple words. We use the part-of-speech tagger to find the part of speech (POS) of every word in the concept description, categorizing them in categories like nouns and verbs. This is used to remove synsets of the wrong POS.

%Data sources:
%WikiPedia
%My application lets you enter the name of a concept. It locates the corresponding WikiPedia article, copying its abstract which is the first paragraph of the article. This text is used as the basis for finding the properties of the concept using natural language processing tools and techniques.

%Node libraries:
%Natural Language Processing tools:
%wordnet-magic
%stanford-simple-nlp?
%includes
%Stanford Parser
%Stanford Tagger
%Pos Tagger

%Javascript libraries:
%Vis.js

%Languages:
%Javascript

%Frameworks:
%AngularJS
%Express

%Environments:

%Node.JS
%Sourcetree
%Git
%GitHub
%TeXstudio
%LaTeX
%Trello

%----------------------------------------------------------------------------------------

\section{Evolution of software}

%----------------------------------------------------------------------------------------

\section{Application structure}
A simplifed overview of the application, its components and data flow can be seen in \textbf{figure 4.1} above.

\begin{figure}
\centering
\includegraphics[width=1\linewidth]{"Figures/Application structure"}
\caption{Overview of the application with the flow of data through the key methods}
\label{fig:application-structure}
\end{figure}

\subsection{Concept}
We have two concepts we want to blend. Each concept has the following data:
- A name which we will use to find the corresponding WikiPedia article for
- A WikiPedia abstract
- The WikiPedia abstract without stopwords
- A list of properties.

\subsection{Wikipedia Abstract}
The WikiPedia abstract contains a text with the description of the concept.

\subsection{Properties}
Each property has the following data:
- The original word
- A list containing the original word and eventual stems found by the stemming algorithm.
- The part-of-speech such as noun, verb or adjective which have been tagged by the POS tagger
- A list of synsets, any unique meaning that the word or the stem of the word may have. Can be pruned to store only synsets of the same part-of-speech that the POS tagger found.
- A reference to the synset of the most likely definition found by the \emph{get most likely definition} method
- A list of hypernyms, from the most specific to the most general, starting with the direct ancestor of the synset.

\subsection{Generic Space}
A list of synsets containing synsets and hypernym synsets that exist in both concepts.

\subsection{Subtrees}
Each concept has a list of subtrees containing unique properties and their hypernyms that are subsuming but are not themselves in the generic space. They are given a compatibility score by the \emph{rank compatible blending suggestions} method. The subtree is assigned a depth, a branch length and a penalty value. The depth value contains its least general hypernym in the generic space.
The branch length is the number of nodes in the subtree. The penalty value is equal to a penalty constant if the property is or has a hypernym in a penalized category.

\subsection{Rank compatible blending suggestions}
Each subtree are given a compatibility score. This is calculated by the depth minus branch length minus penalty. They are then sorted by this compatibility score.

%---------------------------------------------------------------------------------------- 
% Chapter 5

\chapter{Concept Blending Algorithm} % Main chapter title

\label{Chapter5} % For referencing the chapter elsewhere, use \ref{Chapter5} 

%----------------------------------------------------------------------------------------

\section{Concept blending}
My implementation of a concept blending algorithm consists of three steps. The first is retrieving the input data consisting of concept descriptions from a data source such as WikiPedia. The second step is creating the generic space, the structure deriving from both inputs. This is done by retrieving generalizations of the words from each concept, creating a tree structure of shared hypernyms. The third step is finding suggestions of elements to blend by using a heuristic function to rank the elements using patterns based on characteristics we have found in good blending examples.

%----------------------------------------------------------------------------------------

\section{Preparing input data}

\subsection{Finding useful properties}
We have two concepts that we want to blend, for example Cajón and Clapsticks.
We can retrieve the abstracts using a HTTP call to the wikipedia API:

\noindent\fbox{
	\parbox{\textwidth}{
		\url{http://en.wikipedia.org/w/api.php?action=parse&page=conceptName&prop=text&section=0&format=json&redirects=1&callback=?}
	}
}
The parameter ConceptName stands for the name of our concept. By extracting the first <p> paragraph element in the returned JSON string, we can retrieve the raw abstract. For Cajón this abstract is:
 
\noindent\fbox{
	\parbox{\textwidth}{
		A cajón is nominally a box-shaped percussion instrument originally from Peru, played by slapping the front or rear faces with the hands, fingers, or sometimes various implements such as brushes, mallets, or sticks.
	}
}

In order for the algorithm to be able to provide suggestions for the best properties to transfer between two concepts, we need to find properties in abstract that may represent the concept. To be properly able to query WordNet with words from the text we need to process the abstract with the steps seen in \textbf{procedure 1}. Initially the descriptions consists of a lot of common words we are not interested in called stop-words. We use a stop-word list to remove them. But before they are removed we send the original abstract to the Stanford POS tagger to have each word tagged. This data is useful when we want to reduce the number of potential word meanings in the disambiguation step.
The abstract of cajón after processing and removal of stopwords and duplicate words:
\noindent\fbox{
	\parbox{\textwidth}{
		cajón nominally box-shaped percussion instrument originally peru played slapping front rear faces hands fingers sometimes various implements brushes mallets sticks
	}
}
%First the algorithm tries to grammatically tag each word using the Stanford part-of-speech tagger.

\begin{algorithm}
	\caption{Process a WikiPedia abstract to prepare words for WordNet queries}\label{euclid}
	\begin{algorithmic}[1]
		\Require wikipediaAbstract,
		\Ensure processedWikipediaAbstract, wikipediaAbstractWithoutStopwords
		\Procedure{processWikipediaAbstract}{}
		\For {\textit{wikipediaAbstract}}
		\State Remove extra newline
		\State Replace line breaks with whitespaces
		\State Split words separated by forward slash
		\State Remove \&\#160; to avoid the number 160 appearing in the text
		\State Remove [] and the containing text (e.g reference tags)
		\State Remove paranthesis and the containing text
		\State Remove forward-slash and the containing text (e.g phonetics)
		\State Remove - at the end of words
		\State Remove d' at the start of french words
		\State Remove l' at the start of french words
		\EndFor
		\State $\textit{wikipediaAbstractWithoutStopwords} \gets \textit{wikipediaAbstract}$
		\For {\textit{wikipediaAbstractWithoutStopwords}}
		\State Remove 's at the end of words
		\State Remove all apostrophes
		%\State Remove symbols \texttt{[!@#$%^&*(){}?<>+:;",.♭♯£€¤=\`_~}
		\State Remove symbols $!@\#\$\%\&*(){}?<>+\hat{}:;",.'\flat\sharp=\textbackslash\grave{}\_\tilde{}\pounds$\texteuro\textcurrency
		\State Remove words that has numbers in it
		\State Change to lower case (WordNet requires lowercase string)
		
		\State Remove stopwords
		\State Remove duplicate words
		\EndFor
		
		\State \Return [\textit{wikipediaAbstract},  \textit{wikipediaAbstractWithoutStopwords}]
		\EndProcedure
	\end{algorithmic}
\end{algorithm}

%----------------------------------------------------------------------------------------

\subsection{Identifying noun phrases}
When analysing a text for useful properties, some of them may be described by multiple words or a noun phrase. Take for example the two-worded property \emph{sustaining pedal}, which is one of the pedals on a piano that lifts the dampers from the strings to let them continue vibrating. The word sustaining and pedal can not sufficiently describe the sustaining pedal on their own. The word sustaining depends on the word pedal, and it may not be clear if the word pedal is a sustaining pedal, a soft pedal or any other pedal unrelated to pianos. Therefore identifying noun phrases can lead to finding more useful properties in the text. We can identify noun phrases by using the Stanford parser which takes a sentence and label the noun phrases with a NP tag. We can then remove unnecessary parts of the noun phrase such as determinants. If the noun phrase contains more than two words, we can recursively check the tail for other noun phrases. For example in the noun phrase \emph{electric five-string bass} there is also the noun phrase \emph{five-string bass}.

%----------------------------------------------------------------------------------------

\subsection{WordNet database}
WordNet provides a set of synsets for each word in the database, and each of these synsets has a unique definition. If we search for the word \emph{string}, WordNet returns a set of 10 nouns and 7 verbs. We want to search WordNet for every word we found in the abstract and find the correct synsets so that we have a list of unambiguous properties representing the concept. WordNet also store the POS of a word, so we can use the POS tagger data to remove every synset of the wrong POS. For instance if we have tagged \emph{piano} as a noun, we can remove the adjective \emph{piano} which means soft.
%Nevne tagger her? Kan luke vekk nouns eller verbs

\subsection{Stemming}
Sometimes WordNet gives no results because the word is in a different variant than the one in the WordNet database. For instance if the word is \emph{guitars}, it does not match \emph{guitar} in the database. To return \emph{guitar}, we can use a stemming algorithm to find the stem of the word. Often the stem is found in WordNet. This method is not always perfect though. Sometimes the stemmer can over-stem, cutting too much of the word leading to a word that has a different meaning. %Example?
%Pos tagging can avoid this.
%Vurdere lemmatizer istaden då dei gir bedre meining? Updates = updat i stemmer. Lemmatizer bruker wordnet

%----------------------------------------------------------------------------------------

\subsection{Word sense disambiguation}
Now that we have the key words and noun phrases extracted from the abstract, we will try to find the correct meaning of each word. Let's say we want to find the synset corresponding to the \emph{bow} used to play the \emph{cello}. How can we do this? The following example is done using the algorithm in \textbf{procedure 2} shown below. We have to get a list of words that we will look for in each definition of \emph{bow}. We call this list \emph{scoring words}. \emph{Scoring words} will consist of each word in the processed abstract plus a set of musical context words that can act as the musical context if the abstract lack sufficient musical words. Most likely the definitions we are looking for are related to a musical context. The \emph{context words} we use in this example are:
\noindent\fbox{
	\parbox{\textwidth}{music, instrument, tone, tones, sounds, sounding, rhythm, melody, drum, flute, percussion, idiophone, bass, baritone, tenor, alto, soprano, note, pitch, audio, auditory, tune, tuning, acoustic}
}
The \emph{scoring words} combines these with the words of the \emph{cello} abstract:

\noindent\fbox{
	\parbox{\textwidth}{
		cello, violoncello, bowed, plucked, string, instrument, four, strings, tuned, perfect, fifths, low, octave, lower, viola, violin, family, musical, instruments, includes, double, bass, solo, chamber, music, ensembles, orchestras, section, symphony, types, rock, bands, second-largest, lowest, modern, orchestra, largest, pitch
	}
}
We will remove duplicates of words that are in both lists so that each word is equally influential on the scoring. Duplicate words are:

\noindent\fbox{
	\parbox{\textwidth}{music, instrument, bass, pitch}
}

To add score to a definition, we count the number of words (excluding stop-words) in each synset definition which also exists in the \emph{scoring words}. The first definition of \emph{bow} is:

\noindent\fbox{\parbox{\textwidth}{
		a knot with two loops and loose ends; used to tie shoelaces
}}
None of these words are in the \emph{scoring words} of a cello, therefore we give it a score of 0. The next definition is:

\noindent\fbox{\parbox{\textwidth}{
		a slightly curved piece of resilient wood with taut horsehair strands; used in playing certain stringed instruments
}}
One word, instrument, is found in the cello scoring words, therefore we give it a score of 1. We select the synset that has the highest score, which in this case was the correct bow definition. This approach is heuristic and will not always find the correct synset.
Sometimes the word is similar to, but not equal to the scoring word. For example if the definition had the word \emph{bow}, and scoring words had \emph{bowed}, we would not find it. One way to match these would be to search for \emph{bow} in the start of the string \emph{bowed}, and also search for \emph{bowed} in the start of \emph{bow}. Now we would get matches when the scoring word starts with the definition word and vice versa. An alternative way could be to use a stemming algorithm on the words and see if either the scoring word or scoring word stem matches the definition word or definition word stem.

An alternative additional scoring system could be to use \emph{context synsets}, specific synsets that automatically give an extra point for being relevant to the category. For example the musical synset of the word \emph{instrument} would automatically get an extra point over non-musical synsets of instruments like the more general synset \emph{device}.

\begin{algorithm}
	\caption{Get the definition which is most likely to fit the word}\label{euclid}
	\begin{algorithmic}[1]
		\Require ambigousWord,
		definitions,
		scoringWords
		\Ensure highest scoring definition
		\Procedure{GetMostLikelyDefinition}{}
		\For {\textbf{each} \textit{definition} in \textit{definitions}}
		\State $\textit{definitionScore} \gets 0$
		\For {\textbf{each} \textit{definitionWord} in \textit{definition}}
		\For {\textbf{each} \textit{scoringWord} in \textit{scoringWords}}
		\If {\textit{scoringWord} matches \textit{ambigousWord}}
		\State \textbf{continue}
		\EndIf
		\If {\textit{scoringWord} is the first part of \textit{definitionWord}
		\textbf{or} \textit{definitionWord} is the first part of \textit{scoringWord}}
		\State ${\textit{definitionScore} \gets {\textit{definitionScore}+1}}$
		\EndIf
		\EndFor
		\EndFor
		\If {$\textit{definitionScore} > \textit{definitionMaxScore}$}
		\State ${\textit{definitionMaxScore}} \gets {\textit{definitionScore}}$
		\State ${\textit{topDefinition}} \gets {\textit{definition}}$
		\EndIf
		\EndFor
		\If {there is no \textit{topDefinition}}
		\Return null
		\EndIf
		\State \Return \textit{topDefinition}
		\EndProcedure
	\end{algorithmic}
\end{algorithm}

%----------------------------------------------------------------------------------------

\section{Generic space}
\subsection{Generalization using hypernyms}
Our approach to concept blending is inspired by the use of amalgams. When using WordNet we can easily generalize since we have access to each synsets tree of hypernyms. So if we consider the concepts \emph{a red French vehicle} and \emph{a German minivan}, we can find that \emph{French} and \emph{German} share the hypernym \emph{nation}, and minivan and vehicle share the hypernym \emph{vehicle}.

Our solution give score proportional to the speciality of the word and penalize proportionally to how general the generalization is, aiming for the least general generalization. In the example above, minivan should get a good score, since it is specific and the other concept is a vehicle, which makes not a too general generalization. However a minivan and SUV would give a better score, since they are both cars, which is less general than vehicle.

The algorithm we use to create the generic space between is shown below as \textbf{procedure 3}. We can pass \emph{Cajón} as the source concept and \emph{Clapsticks} as the target concept. After processing the abstracts and finding the most likely synsets of each property, we can use WordNet to retrieve a list of every hypernym ancestor of each synset, which will go from the least general to the most general generalization. They can be accessed from the array \emph{concept.property.hypernyms}. To create the generic space, we simply retrieve all synsets and hypernymsynsets in each concept and use a intersection set operation to keep what we find in both concepts.

\begin{algorithm}
	\caption{Get all synsets in the generic space between two concepts}\label{euclid}
	\begin{algorithmic}[1]
		\Require sourceConcept,
		targetConcept
		\Ensure highest scoring definition
		\Procedure{GetGenericSpace}{}
		\For {\textbf{each} \textit{property} in \textit{sourceConcept.properties}}
		\State ${\textit{sourceSynsets} \gets {\textit{sourceSynsets} + \textit{property}}}$
		\For {\textbf{each} \textit{hypernym} in \textit{property.hypernyms}}
		\State ${\textit{sourceSynsets} \gets {\textit{sourceSynsets} + \textit{hypernym}}}$
		\EndFor
		\EndFor
		
		\For {\textbf{each} \textit{property} in \textit{targetConcept.properties}}
		\State ${\textit{targetSynsets} \gets {\textit{targetSynsets} + \textit{property}}}$
		\For {\textbf{each} \textit{hypernym} in \textit{property.hypernyms}}
		\State ${\textit{targetSynsets} \gets {\textit{targetSynsets} + \textit{hypernym}}}$
		\EndFor
		\EndFor
		\State ${\textit{sourceSynsets} \gets {removeDuplicates(\textit{sourceSynsets})}}$
		\State ${\textit{targetSynsets} \gets {removeDuplicates(\textit{targetSynsets})}}$
		\State \Return ${SetOperations.intersection(\textit{sourceSynsets}, \textit{targetSynsets})}$
		\EndProcedure
	\end{algorithmic}
\end{algorithm}

\subsection{Finding matches}
In order to provide good suggestions, we need to find properties that share a hypernym in the generic space with some property in the other concept. For instance a football field and a golf terrain share a common hypernym \emph{piece of land} which means they show a kind of commonality and could therefore be more compatible for blending than less related pairs. 

%----------------------------------------------------------------------------------------

\section{Finding good blending suggestions}
We came up with a solution where properties from each concept are sorted with a score that represents how well they are suited for blending. We were improving the algorithm by identifying patterns typical of good candidates in the tree structure of the property words and their hypernyms.
\subsection{Depth of the word in hypernym tree}
The first pattern we found was that words that have high generality are not useful. We want more specific words like \emph{drumstick} rather than general and vague words like \emph{object}. Therefore we gave points equal to the depth of the word in the tree.
\subsection{Branch length to shared hypernyms}
The second pattern we found was when properties of two concepts share a common root or hypernym, the length of the branch from this root seemed to be useful information. The longer the branch length, the worse were the candidate. Therefore we punished the word by subtracting points equal to the branch length.
\subsection{Penalizing certain categories}
Now the winning words were specific, but too often abstract, being an instrument in itself or a person (e.g. the inventor). Therefore we chose to punish the words significantly when being in the branch of \emph{abstract entities}, \emph{instruments} or \emph{individuals}. The resulting suggestions were then both physical and specific, which seems to be optimal for creating new instruments.
\subsection{Removing properties in penalized categories}
An observation of the properties that connected in the general space, showed that a lot of these were in penalized categories. For example a property of the penalized category instrument is a device, which makes devices in the other concept score higher than if there were no other devices in the first concept. A way to avoid this could be to remove properties of penalized categories altogether. This leads to far less blending suggestions, but the suggestions that are left are of relatively high quality. Since no blending suggestions are in penalized categories also makes the penalty score redundant.

\subsection{Ranking Compatible Blending Suggestions}
First we find each property that are not in the generic space themselves but has a hypernym in the generic space. This means that this property can be blended and we can give it a compatibility score for whether it is a good blending suggestion. We call the branch from the synset and hypernyms subsuming its hypernyms in the generic space for \emph{blending subtree}. We set the compatibility score to the depth of its least general hypernym in the generic space minus the length of the \emph{blending subtree}. Finally we can optionally subtract a constant penalty value if the property is or has a hypernym in a penalized category.  With this score we can store the unique properties in a list which we can rank using a heuristic function. The algorithm for finding unique properties and ranking them is shown below in \textbf{procedure 4}.

\begin{algorithm}
	\caption{Perform creative strategy by ranking a concepts subtrees by finding a pair in the other concept with the highest compability, defined by the expression Depth - BranchLength - Penalty}\label{euclid}
	\begin{algorithmic}[1]
		\Require properties,
		concept,
		genericSpace,
		penaltyCategories,
		penaltyPoints
		\Ensure a list of branches ranked by their blending compatibility score
		\Procedure{RankCompatibleBlendingSuggestions}{}
		\For {\textbf{each} \textit{property} in \textit{properties}}
		\If {$\textit{property.name} = \textit{concept.name}$
		\\\qquad\qquad\textbf{or} \textit{property} has not been disambiguated
		\\\qquad\qquad\textbf{or} \textit{property} is in \textit{genericSpace}}
		\State \textbf{continue}
		\EndIf
		\For {\textit{i} = 0 to \textit{property.hypernyms.length}}
		\If {\textit{property.hypernyms[i]} is not in \textit{genericSpace}}
		\State {\textbf{continue}}
		\EndIf
		\State ${\textit{subtree.branch} \gets {[\textit{property}] \text{ concatinated with } \textit{property.hypernyms[0:i]}}}$
		\State ${\textit{depth} \gets {\textit{property.hypernyms.length} - i}}$
		\State ${\textit{branchLength} \gets {\textit{subtree.branch.length}}}$
		\If {\textit{subtree.branch} has an element in \textit{penaltyCategories}}
		\State ${\textit{penalty} \gets \textit{penaltyPoints}}$
		\EndIf
		\State ${\textit{subtree.compatibility} \gets {\textit{depth} - \textit{branchLength} - \textit{penalty}}}$
		\State ${\textit{subtrees} \gets {\textit{subtrees} \text{ concatinated with } \textit{subtree}}}$
		\State {\textbf{break} inner for loop}
		\EndFor
		\EndFor
		\State \Return \textit{subtrees} \textbf{sorted} by \textit{subtree.compability} in descending order
	\EndProcedure
\end{algorithmic}
\end{algorithm}

%----------------------------------------------------------------------------------------

\section{Future work} 
% Chapter 6

\chapter{Results} % Main chapter title

\label{Chapter6} % For referencing the chapter elsewhere, use \ref{Chapter6} 

%----------------------------------------------------------------------------------------

\section{Musical instruments}
When blending musical instruments with no penalty categories, we got the top blending suggestions seen in \textbf{table 6.2}. Top synsets such as \emph{five} and \emph{running} where in the category \emph{abstract entities} and on their own not very interesting in describing musical instruments. If they were to be useful, we would need to model their original relation in the context, which was \emph{five strings} and \emph{running along the soundboard}.
When blending musical instruments using the penalty score 10 for the categories abstract entity, instrument and individual, we got the top blending suggestions seen in \textbf{table 6.2}. We did not exclude synsets that were in the penalized categories. This meant that top synsets often got a good score although the synset in the other instrument that shared the same hypernym was in a penalized category.
\begin{table}
	\caption{Examples of best suggestions when blending different pairs of instruments. Score is based on depth of word in tree (D) minus branch length from shared hypernym (BL)}
	\label{tab:topscoringexamples}
	\centering
	\begin{tabular}{l p{30mm} p{30mm} l l}
		\toprule
		\tabhead{Instrument} & \tabhead{Top scoring branch} & \tabhead{Context} & \tabhead{D - BL} & \tabhead{Score} \\
		\midrule
		Banjo &
		integer \newline
		digit \newline
		\emph{five} & \emph{five} or six-stringed instrument with a thin membrane stretched over a frame or cavity as a resonator & 7 - 3 & 4	\\
		Melodica &
		device \newline 
		\emph{keyboard} &
		The \emph{keyboard} is usually two or three octaves long & 8 - 2 &	6
		\\
		\midrule
		Harp &
		idea \newline
		plan \newline
		plan of action \newline
		play \newline
		football play \newline
		\emph{running} \newline
		(wrong synset) & The harp is a stringed musical instrument that has a number of individual strings \emph{running} at an angle to its soundboard; the strings are plucked with the fingers & 7 - 6 & 1	\\
		 &
		construction \newline
		area \newline
		enclosure \newline
		chamber \newline
		cavity resonator \newline
		\emph{soundboard} &
		The harp is a stringed musical instrument that has a number of individual strings running at an angle to its \emph{soundboard}; the strings are plucked with the fingers & 7 - 6 &	1 \\
		Violin &
		instrumentalist \newline
		\emph{violinists} & \emph{Violinists} and collectors particularly prize the instruments made by the Stradivari & 11 - 2 & 9	\\
		\\
		\bottomrule\\
	\end{tabular}
\end{table}

\begin{table}
	\caption{Examples of best suggestions when blending different pairs of instruments. Score is based on depth of word in tree (D) minus branch length from shared hypernym (BL). Words are also penalized if they have a hypernym called \emph{abstract entity}, \emph{instrument} or \emph{individual}}
	\label{tab:topscoringexamplespenalty}
	\centering
	\begin{tabular}{l p{30mm} p{30mm} l l}
		\toprule
		\tabhead{Instrument} & \tabhead{Top scoring branch} & \tabhead{Context} & \tabhead{D - BL} & \tabhead{Score} \\
		\midrule
		Banjo &
		artefact \newline
		flat solid \newline
		\emph{membrane} & five or six-stringed instrument with a thin \emph{membrane} stretched over a frame or cavity as a resonator & 6 - 3 & 3	\\
		&
		body part \newline
		system \newline
		\emph{frame} \newline
		(wrong synset)
		& five or six-stringed instrument with a thin membrane stretched over a \emph{frame} or cavity as a resonator & 6 - 3 & 3	\\
		Melodica &
		device \newline 
		\emph{keyboard} &
		The \emph{keyboard} is usually two or three octaves long & 8 - 2 &	6
		\\
		\midrule
		Harp &
		construction \newline
		area \newline
		enclosure \newline
		chamber \newline
		cavity resonator \newline
		\emph{soundboard} & The harp is a stringed musical instrument that has a number of individual strings running at an angle to its \emph{soundboard}; the strings are plucked with the fingers & 7 - 6 & 1	\\
		Violin &
		instrumentality \newline
		implement \newline
		stick \newline
		\emph{bow} & and is most commonly played by drawing a \emph{bow} across its strings \newline & 7 - 4 & 3	\\
		&
		construction \newline
		building \newline
		house \newline
		\emph{cottage} & 
		as well as still greater numbers of mass-produced commercial "trade violins" coming from \emph{cottage} industries in places such as saxony & 7 - 4 & 3	
		\\
		\bottomrule\\
	\end{tabular}
\end{table}

%----------------------------------------------------------------------------------------

\section{Alternative blending category}
\subsection{Sports}
The algorithm was tested on the alternative blending category sports. The context words to help the scoring of synsets were: 'sport', 'athletic', 'ball', 'field', 'team', 'goal', 'rules', 'foul', 'score', 'player', 'referee', 'match', 'racket', 'club', 'net', 'skates', 'skis', 'helmet', 'pads', 'bat', 'pitch', 'court', 'tee', 'green', 'win', 'draw' and 'loss'.
The penalized category was 'athletic game', as we would not want sports to be suggested, since it is the category itself.

When performing the algorithm on \emph{soccer} and \emph{golf}, it suggested to move \emph{eleven} (squad) from \emph{soccer} to \emph{golf}, with a depth of 7 minus branch length 1, giving a score of 6. It suggested to move \emph{terrains} from \emph{golf} to \emph{soccer}, with a depth of 7 minus branch length 1, giving a score of 6. The results are interesting, since it can be interpreted as suggesting golf as a team sport of eleven players. Although golf is usually played as an individual sport, there are also a number of team tournaments. Playing soccer on uneven terrain is also an interesting idea that is not common but would be possible to organize, although it may not be an evenly balanced game. One reason for why the interpretations of these suggestions worked well may be because the relations to the other elements in the concept did not contradict. \emph{Terrains} in golf matched with the part \emph{field} in soccer through their common hypernym \emph{parcel}. By not going further up the generalizing chain, the more likely relations still work.

%----------------------------------------------------------------------------------------

\section{Evaluation of blending}

%---------------------------------------------------------------------------------------- 
%% Chapter 7

\chapter{Evaluation} % Main chapter title

\label{Chapter7} % For referencing the chapter elsewhere, use \ref{Chapter7} 

%----------------------------------------------------------------------------------------

\section{Performance/Accuracy}

%----------------------------------------------------------------------------------------

\section{Creative strategy}

%----------------------------------------------------------------------------------------

\section{Creative value}

%----------------------------------------------------------------------------------------

\section{Improvements}

%---------------------------------------------------------------------------------------- 
% Chapter 8

\chapter{Discussion} % Main chapter title

\label{Chapter8} % For referencing the chapter elsewhere, use \ref{Chapter8} 

%----------------------------------------------------------------------------------------

\section{Theory}

%----------------------------------------------------------------------------------------

\section{User input in the algorithm}

%----------------------------------------------------------------------------------------

\section{User interpretation of blending results}

%----------------------------------------------------------------------------------------

\section{Different uses of the concept blending implementation}

%----------------------------------------------------------------------------------------

\section{Obstacles}

%----------------------------------------------------------------------------------------

\section{Summary}

%---------------------------------------------------------------------------------------- 
% Chapter 9

\chapter{Conclusion and future work} % Main chapter title

\label{Chapter9} % For referencing the chapter elsewhere, use \ref{Chapter9} 

%----------------------------------------------------------------------------------------

\section{Conclusions}

%----------------------------------------------------------------------------------------

\section{Future work}

%----------------------------------------------------------------------------------------

\section{Future research}

%---------------------------------------------------------------------------------------- 
%% Chapter 10

\chapter{LaTeX Guide} % Main chapter title

\label{Chapter10} % For referencing the chapter elsewhere, use \ref{Chapter10} 

%----------------------------------------------------------------------------------------

\section{Welcome and Thank You}
Welcome to this \LaTeX{} Thesis Template, a beautiful and easy to use template for writing a thesis using the \LaTeX{} typesetting system.

If you are writing a thesis (or will be in the future) and its subject is technical or mathematical (though it doesn't have to be), then creating it in \LaTeX{} is highly recommended as a way to make sure you can just get down to the essential writing without having to worry over formatting or wasting time arguing with your word processor.

\LaTeX{} is easily able to professionally typeset documents that run to hundreds or thousands of pages long. With simple mark-up commands, it automatically sets out the table of contents, margins, page headers and footers and keeps the formatting consistent and beautiful. One of its main strengths is the way it can easily typeset mathematics, even \emph{heavy} mathematics. Even if those equations are the most horribly twisted and most difficult mathematical problems that can only be solved on a super-computer, you can at least count on \LaTeX{} to make them look stunning.

%----------------------------------------------------------------------------------------

\section{Learning \LaTeX{}}

\LaTeX{} is not a \textsc{wysiwyg} (What You See is What You Get) program, unlike word processors such as Microsoft Word or Apple's Pages. Instead, a document written for \LaTeX{} is actually a simple, plain text file that contains \emph{no formatting}. You tell \LaTeX{} how you want the formatting in the finished document by writing in simple commands amongst the text, for example, if I want to use \emph{italic text for emphasis}, I write the \verb|\emph{text}| command and put the text I want in italics in between the curly braces. This means that \LaTeX{} is a \enquote{mark-up} language, very much like HTML.

\subsection{A (not so short) Introduction to \LaTeX{}}

If you are new to \LaTeX{}, there is a very good eBook -- freely available online as a PDF file -- called, \enquote{The Not So Short Introduction to \LaTeX{}}. The book's title is typically shortened to just \emph{lshort}. You can download the latest version (as it is occasionally updated) from here:
\url{http://www.ctan.org/tex-archive/info/lshort/english/lshort.pdf}

It is also available in several other languages. Find yours from the list on this page: \url{http://www.ctan.org/tex-archive/info/lshort/}

It is recommended to take a little time out to learn how to use \LaTeX{} by creating several, small `test' documents, or having a close look at several templates on:\\ 
\url{http://www.LaTeXTemplates.com}\\ 
Making the effort now means you're not stuck learning the system when what you \emph{really} need to be doing is writing your thesis.

\subsection{A Short Math Guide for \LaTeX{}}

If you are writing a technical or mathematical thesis, then you may want to read the document by the AMS (American Mathematical Society) called, \enquote{A Short Math Guide for \LaTeX{}}. It can be found online here:
\url{http://www.ams.org/tex/amslatex.html}
under the \enquote{Additional Documentation} section towards the bottom of the page.

\subsection{Common \LaTeX{} Math Symbols}
There are a multitude of mathematical symbols available for \LaTeX{} and it would take a great effort to learn the commands for them all. The most common ones you are likely to use are shown on this page:
\url{http://www.sunilpatel.co.uk/latex-type/latex-math-symbols/}

You can use this page as a reference or crib sheet, the symbols are rendered as large, high quality images so you can quickly find the \LaTeX{} command for the symbol you need.

\subsection{\LaTeX{} on a Mac}
 
The \LaTeX{} distribution is available for many systems including Windows, Linux and Mac OS X. The package for OS X is called MacTeX and it contains all the applications you need -- bundled together and pre-customized -- for a fully working \LaTeX{} environment and work flow.
 
MacTeX includes a custom dedicated \LaTeX{} editor called TeXShop for writing your `\file{.tex}' files and BibDesk: a program to manage your references and create your bibliography section just as easily as managing songs and creating playlists in iTunes.

%----------------------------------------------------------------------------------------

\section{Getting Started with this Template}

If you are familiar with \LaTeX{}, then you should explore the directory structure of the template and then proceed to place your own information into the \emph{THESIS INFORMATION} block of the \file{main.tex} file. You can then modify the rest of this file to your unique specifications based on your degree/university. Section \ref{FillingFile} on page \pageref{FillingFile} will help you do this. Make sure you also read section \ref{ThesisConventions} about thesis conventions to get the most out of this template.

If you are new to \LaTeX{} it is recommended that you carry on reading through the rest of the information in this document.

Before you begin using this template you should ensure that its style complies with the thesis style guidelines imposed by your institution. In most cases this template style and layout will be suitable. If it is not, it may only require a small change to bring the template in line with your institution's recommendations. These modifications will need to be done on the \file{MastersDooralThesis.cls} file.

\subsection{About this Template}

This \LaTeX{} Thesis Template is originally based and created around a \LaTeX{} style file created by Steve R.\ Gunn from the University of Southampton (UK), department of Electronics and Computer Science. You can find his original thesis style file at his site, here:
\url{http://www.ecs.soton.ac.uk/~srg/softwaretools/document/templates/}

Steve's \file{ecsthesis.cls} was then taken by Sunil Patel who modified it by creating a skeleton framework and folder structure to place the thesis files in. The resulting template can be found on Sunil's site here:
\url{http://www.sunilpatel.co.uk/thesis-template}

Sunil's template was made available through \url{http://www.LaTeXTemplates.com} where it was modified many times based on user requests and questions. Version 2.0 and onwards of this template represents a major modification to Sunil's template and is, in fact, hardly recognisable. The work to make version 2.0 possible was carried out by \href{mailto:vel@latextemplates.com}{Vel} and Johannes Böttcher.

%----------------------------------------------------------------------------------------

\section{What this Template Includes}

\subsection{Folders}

This template comes as a single zip file that expands out to several files and folders. The folder names are mostly self-explanatory:

\keyword{Appendices} -- this is the folder where you put the appendices. Each appendix should go into its own separate \file{.tex} file. An example and template are included in the directory.

\keyword{Chapters} -- this is the folder where you put the thesis chapters. A thesis usually has about six chapters, though there is no hard rule on this. Each chapter should go in its own separate \file{.tex} file and they can be split as:
\begin{itemize}
\item Chapter 1: Introduction to the thesis topic
\item Chapter 2: Background information and theory
\item Chapter 3: (Laboratory) experimental setup
\item Chapter 4: Details of experiment 1
\item Chapter 5: Details of experiment 2
\item Chapter 6: Discussion of the experimental results
\item Chapter 7: Conclusion and future directions
\end{itemize}
This chapter layout is specialised for the experimental sciences, your discipline may be different.

\keyword{Figures} -- this folder contains all figures for the thesis. These are the final images that will go into the thesis document.

\subsection{Files}

Included are also several files, most of them are plain text and you can see their contents in a text editor. After initial compilation, you will see that more auxiliary files are created by \LaTeX{} or BibTeX and which you don't need to delete or worry about:

\keyword{example.bib} -- this is an important file that contains all the bibliographic information and references that you will be citing in the thesis for use with BibTeX. You can write it manually, but there are reference manager programs available that will create and manage it for you. Bibliographies in \LaTeX{} are a large subject and you may need to read about BibTeX before starting with this. Many modern reference managers will allow you to export your references in BibTeX format which greatly eases the amount of work you have to do.

\keyword{MastersDoctoralThesis.cls} -- this is an important file. It is the class file that tells \LaTeX{} how to format the thesis. 

\keyword{main.pdf} -- this is your beautifully typeset thesis (in the PDF file format) created by \LaTeX{}. It is supplied in the PDF with the template and after you compile the template you should get an identical version.

\keyword{main.tex} -- this is an important file. This is the file that you tell \LaTeX{} to compile to produce your thesis as a PDF file. It contains the framework and constructs that tell \LaTeX{} how to layout the thesis. It is heavily commented so you can read exactly what each line of code does and why it is there. After you put your own information into the \emph{THESIS INFORMATION} block -- you have now started your thesis!

Files that are \emph{not} included, but are created by \LaTeX{} as auxiliary files include:

\keyword{main.aux} -- this is an auxiliary file generated by \LaTeX{}, if it is deleted \LaTeX{} simply regenerates it when you run the main \file{.tex} file.

\keyword{main.bbl} -- this is an auxiliary file generated by BibTeX, if it is deleted, BibTeX simply regenerates it when you run the \file{main.aux} file. Whereas the \file{.bib} file contains all the references you have, this \file{.bbl} file contains the references you have actually cited in the thesis and is used to build the bibliography section of the thesis.

\keyword{main.blg} -- this is an auxiliary file generated by BibTeX, if it is deleted BibTeX simply regenerates it when you run the main \file{.aux} file.

\keyword{main.lof} -- this is an auxiliary file generated by \LaTeX{}, if it is deleted \LaTeX{} simply regenerates it when you run the main \file{.tex} file. It tells \LaTeX{} how to build the \emph{List of Figures} section.

\keyword{main.log} -- this is an auxiliary file generated by \LaTeX{}, if it is deleted \LaTeX{} simply regenerates it when you run the main \file{.tex} file. It contains messages from \LaTeX{}, if you receive errors and warnings from \LaTeX{}, they will be in this \file{.log} file.

\keyword{main.lot} -- this is an auxiliary file generated by \LaTeX{}, if it is deleted \LaTeX{} simply regenerates it when you run the main \file{.tex} file. It tells \LaTeX{} how to build the \emph{List of Tables} section.

\keyword{main.out} -- this is an auxiliary file generated by \LaTeX{}, if it is deleted \LaTeX{} simply regenerates it when you run the main \file{.tex} file.

So from this long list, only the files with the \file{.bib}, \file{.cls} and \file{.tex} extensions are the most important ones. The other auxiliary files can be ignored or deleted as \LaTeX{} and BibTeX will regenerate them.

%----------------------------------------------------------------------------------------

\section{Filling in Your Information in the \file{main.tex} File}\label{FillingFile}

You will need to personalise the thesis template and make it your own by filling in your own information. This is done by editing the \file{main.tex} file in a text editor or your favourite LaTeX environment.

Open the file and scroll down to the third large block titled \emph{THESIS INFORMATION} where you can see the entries for \emph{University Name}, \emph{Department Name}, etc \ldots

Fill out the information about yourself, your group and institution. You can also insert web links, if you do, make sure you use the full URL, including the \code{http://} for this. If you don't want these to be linked, simply remove the \verb|\href{url}{name}| and only leave the name.

When you have done this, save the file and recompile \code{main.tex}. All the information you filled in should now be in the PDF, complete with web links. You can now begin your thesis proper!

%----------------------------------------------------------------------------------------

\section{The \code{main.tex} File Explained}

The \file{main.tex} file contains the structure of the thesis. There are plenty of written comments that explain what pages, sections and formatting the \LaTeX{} code is creating. Each major document element is divided into commented blocks with titles in all capitals to make it obvious what the following bit of code is doing. Initially there seems to be a lot of \LaTeX{} code, but this is all formatting, and it has all been taken care of so you don't have to do it.

Begin by checking that your information on the title page is correct. For the thesis declaration, your institution may insist on something different than the text given. If this is the case, just replace what you see with what is required in the \emph{DECLARATION PAGE} block.

Then comes a page which contains a funny quote. You can put your own, or quote your favourite scientist, author, person, and so on. Make sure to put the name of the person who you took the quote from.

Following this is the abstract page which summarises your work in a condensed way and can almost be used as a standalone document to describe what you have done. The text you write will cause the heading to move up so don't worry about running out of space.

Next come the acknowledgements. On this page, write about all the people who you wish to thank (not forgetting parents, partners and your advisor/supervisor).

The contents pages, list of figures and tables are all taken care of for you and do not need to be manually created or edited. The next set of pages are more likely to be optional and can be deleted since they are for a more technical thesis: insert a list of abbreviations you have used in the thesis, then a list of the physical constants and numbers you refer to and finally, a list of mathematical symbols used in any formulae. Making the effort to fill these tables means the reader has a one-stop place to refer to instead of searching the internet and references to try and find out what you meant by certain abbreviations or symbols.

The list of symbols is split into the Roman and Greek alphabets. Whereas the abbreviations and symbols ought to be listed in alphabetical order (and this is \emph{not} done automatically for you) the list of physical constants should be grouped into similar themes.

The next page contains a one line dedication. Who will you dedicate your thesis to?

Finally, there is the block where the chapters are included. Uncomment the lines (delete the \code{\%} character) as you write the chapters. Each chapter should be written in its own file and put into the \emph{Chapters} folder and named \file{Chapter1}, \file{Chapter2}, etc\ldots Similarly for the appendices, uncomment the lines as you need them. Each appendix should go into its own file and placed in the \emph{Appendices} folder.

After the preamble, chapters and appendices finally comes the bibliography. The bibliography style (called \option{authoryear}) is used for the bibliography and is a fully featured style that will even include links to where the referenced paper can be found online. Do not underestimate how grateful your reader will be to find that a reference to a paper is just a click away. Of course, this relies on you putting the URL information into the BibTeX file in the first place.

%----------------------------------------------------------------------------------------

\section{Thesis Features and Conventions}\label{ThesisConventions}

To get the best out of this template, there are a few conventions that you may want to follow.

One of the most important (and most difficult) things to keep track of in such a long document as a thesis is consistency. Using certain conventions and ways of doing things (such as using a Todo list) makes the job easier. Of course, all of these are optional and you can adopt your own method.

\subsection{Printing Format}

This thesis template is designed for double sided printing (i.e. content on the front and back of pages) as most theses are printed and bound this way. Switching to one sided printing is as simple as uncommenting the \option{oneside} option of the \code{documentclass} command at the top of the \file{main.tex} file. You may then wish to adjust the margins to suit specifications from your institution.

The headers for the pages contain the page number on the outer side (so it is easy to flick through to the page you want) and the chapter name on the inner side.

The text is set to 11 point by default with single line spacing, again, you can tune the text size and spacing should you want or need to using the options at the very start of \file{main.tex}. The spacing can be changed similarly by replacing the \option{singlespacing} with \option{onehalfspacing} or \option{doublespacing}.

\subsection{Using US Letter Paper}

The paper size used in the template is A4, which is the standard size in Europe. If you are using this thesis template elsewhere and particularly in the United States, then you may have to change the A4 paper size to the US Letter size. This can be done in the margins settings section in \file{main.tex}.

Due to the differences in the paper size, the resulting margins may be different to what you like or require (as it is common for institutions to dictate certain margin sizes). If this is the case, then the margin sizes can be tweaked by modifying the values in the same block as where you set the paper size. Now your document should be set up for US Letter paper size with suitable margins.

\subsection{References}

The \code{biblatex} package is used to format the bibliography and inserts references such as this one \parencite{Reference1}. The options used in the \file{main.tex} file mean that the in-text citations of references are formatted with the author(s) listed with the date of the publication. Multiple references are separated by semicolons (e.g. \parencite{Reference2, Reference1}) and references with more than three authors only show the first author with \emph{et al.} indicating there are more authors (e.g. \parencite{Reference3}). This is done automatically for you. To see how you use references, have a look at the \file{Chapter1.tex} source file. Many reference managers allow you to simply drag the reference into the document as you type.

Scientific references should come \emph{before} the punctuation mark if there is one (such as a comma or period). The same goes for footnotes\footnote{Such as this footnote, here down at the bottom of the page.}. You can change this but the most important thing is to keep the convention consistent throughout the thesis. Footnotes themselves should be full, descriptive sentences (beginning with a capital letter and ending with a full stop). The APA6 states: \enquote{Footnote numbers should be superscripted, [...], following any punctuation mark except a dash.} The Chicago manual of style states: \enquote{A note number should be placed at the end of a sentence or clause. The number follows any punctuation mark except the dash, which it precedes. It follows a closing parenthesis.}

The bibliography is typeset with references listed in alphabetical order by the first author's last name. This is similar to the APA referencing style. To see how \LaTeX{} typesets the bibliography, have a look at the very end of this document (or just click on the reference number links in in-text citations).

\subsubsection{A Note on bibtex}

The bibtex backend used in the template by default does not correctly handle unicode character encoding (i.e. "international" characters). You may see a warning about this in the compilation log and, if your references contain unicode characters, they may not show up correctly or at all. The solution to this is to use the biber backend instead of the outdated bibtex backend. This is done by finding this in \file{main.tex}: \option{backend=bibtex} and changing it to \option{backend=biber}. You will then need to delete all auxiliary BibTeX files and navigate to the template directory in your terminal (command prompt). Once there, simply type \code{biber main} and biber will compile your bibliography. You can then compile \file{main.tex} as normal and your bibliography will be updated. An alternative is to set up your LaTeX editor to compile with biber instead of bibtex, see \href{http://tex.stackexchange.com/questions/154751/biblatex-with-biber-configuring-my-editor-to-avoid-undefined-citations/}{here} for how to do this for various editors.

\subsection{Tables}

Tables are an important way of displaying your results, below is an example table which was generated with this code:

{\small
\begin{verbatim}
\begin{table}
\caption{The effects of treatments X and Y on the four groups studied.}
\label{tab:treatments}
\centering
\begin{tabular}{l l l}
\toprule
\tabhead{Groups} & \tabhead{Treatment X} & \tabhead{Treatment Y} \\
\midrule
1 & 0.2 & 0.8\\
2 & 0.17 & 0.7\\
3 & 0.24 & 0.75\\
4 & 0.68 & 0.3\\
\bottomrule\\
\end{tabular}
\end{table}
\end{verbatim}
}

\begin{table}
\caption{The effects of treatments X and Y on the four groups studied.}
\label{tab:treatments}
\centering
\begin{tabular}{l l l}
\toprule
\tabhead{Groups} & \tabhead{Treatment X} & \tabhead{Treatment Y} \\
\midrule
1 & 0.2 & 0.8\\
2 & 0.17 & 0.7\\
3 & 0.24 & 0.75\\
4 & 0.68 & 0.3\\
\bottomrule\\
\end{tabular}
\end{table}

You can reference tables with \verb|\ref{<label>}| where the label is defined within the table environment. See \file{Chapter1.tex} for an example of the label and citation (e.g. Table~\ref{tab:treatments}).

\subsection{Figures}

There will hopefully be many figures in your thesis (that should be placed in the \emph{Figures} folder). The way to insert figures into your thesis is to use a code template like this:
\begin{verbatim}
\begin{figure}
\centering
\includegraphics{Figures/Electron}
\decoRule
\caption[An Electron]{An electron (artist's impression).}
\label{fig:Electron}
\end{figure}
\end{verbatim}
Also look in the source file. Putting this code into the source file produces the picture of the electron that you can see in the figure below.

\begin{figure}[th]
\centering
\includegraphics{Figures/Electron}
\decoRule
\caption[An Electron]{An electron (artist's impression).}
\label{fig:Electron}
\end{figure}

Sometimes figures don't always appear where you write them in the source. The placement depends on how much space there is on the page for the figure. Sometimes there is not enough room to fit a figure directly where it should go (in relation to the text) and so \LaTeX{} puts it at the top of the next page. Positioning figures is the job of \LaTeX{} and so you should only worry about making them look good!

Figures usually should have captions just in case you need to refer to them (such as in Figure~\ref{fig:Electron}). The \verb|\caption| command contains two parts, the first part, inside the square brackets is the title that will appear in the \emph{List of Figures}, and so should be short. The second part in the curly brackets should contain the longer and more descriptive caption text.

The \verb|\decoRule| command is optional and simply puts an aesthetic horizontal line below the image. If you do this for one image, do it for all of them.

\LaTeX{} is capable of using images in pdf, jpg and png format.

\subsection{Typesetting mathematics}

If your thesis is going to contain heavy mathematical content, be sure that \LaTeX{} will make it look beautiful, even though it won't be able to solve the equations for you.

The \enquote{Not So Short Introduction to \LaTeX} (available on \href{http://www.ctan.org/tex-archive/info/lshort/english/lshort.pdf}{CTAN}) should tell you everything you need to know for most cases of typesetting mathematics. If you need more information, a much more thorough mathematical guide is available from the AMS called, \enquote{A Short Math Guide to \LaTeX} and can be downloaded from:
\url{ftp://ftp.ams.org/pub/tex/doc/amsmath/short-math-guide.pdf}

There are many different \LaTeX{} symbols to remember, luckily you can find the most common symbols in \href{http://ctan.org/pkg/comprehensive}{The Comprehensive \LaTeX~Symbol List}.

You can write an equation, which is automatically given an equation number by \LaTeX{} like this:
\begin{verbatim}
\begin{equation}
E = mc^{2}
\label{eqn:Einstein}
\end{equation}
\end{verbatim}

This will produce Einstein's famous energy-matter equivalence equation:
\begin{equation}
E = mc^{2}
\label{eqn:Einstein}
\end{equation}

All equations you write (which are not in the middle of paragraph text) are automatically given equation numbers by \LaTeX{}. If you don't want a particular equation numbered, use the unnumbered form:
\begin{verbatim}
\[ a^{2}=4 \]
\end{verbatim}

%----------------------------------------------------------------------------------------

\section{Sectioning and Subsectioning}

You should break your thesis up into nice, bite-sized sections and subsections. \LaTeX{} automatically builds a table of Contents by looking at all the \verb|\chapter{}|, \verb|\section{}|  and \verb|\subsection{}| commands you write in the source.

The Table of Contents should only list the sections to three (3) levels. A \verb|chapter{}| is level zero (0). A \verb|\section{}| is level one (1) and so a \verb|\subsection{}| is level two (2). In your thesis it is likely that you will even use a \verb|subsubsection{}|, which is level three (3). The depth to which the Table of Contents is formatted is set within \file{MastersDoctoralThesis.cls}. If you need this changed, you can do it in \file{main.tex}.

%----------------------------------------------------------------------------------------

\section{In Closing}

You have reached the end of this mini-guide. You can now rename or overwrite this pdf file and begin writing your own \file{Chapter1.tex} and the rest of your thesis. The easy work of setting up the structure and framework has been taken care of for you. It's now your job to fill it out!

Good luck and have lots of fun!

\begin{flushright}
Guide written by ---\\
Sunil Patel: \href{http://www.sunilpatel.co.uk}{www.sunilpatel.co.uk}\\
Vel: \href{http://www.LaTeXTemplates.com}{LaTeXTemplates.com}
\end{flushright}
 

%----------------------------------------------------------------------------------------
%	THESIS CONTENT - APPENDICES
%----------------------------------------------------------------------------------------

\appendix % Cue to tell LaTeX that the following "chapters" are Appendices

% Include the appendices of the thesis as separate files from the Appendices folder
% Uncomment the lines as you write the Appendices

%% Appendix A

\chapter{Frequently Asked Questions} % Main appendix title

\label{AppendixA} % For referencing this appendix elsewhere, use \ref{AppendixA}

\section{How do I change the colors of links?}

The color of links can be changed to your liking using:

{\small\verb!\hypersetup{urlcolor=red}!}, or

{\small\verb!\hypersetup{citecolor=green}!}, or

{\small\verb!\hypersetup{allcolor=blue}!}.

\noindent If you want to completely hide the links, you can use:

{\small\verb!\hypersetup{allcolors=.}!}, or even better: 

{\small\verb!\hypersetup{hidelinks}!}.

\noindent If you want to have obvious links in the PDF but not the printed text, use:

{\small\verb!\hypersetup{colorlinks=false}!}.

%\include{Appendices/AppendixB}
%\include{Appendices/AppendixC}

%----------------------------------------------------------------------------------------
%	BIBLIOGRAPHY
%----------------------------------------------------------------------------------------

\printbibliography[heading=bibintoc]

%----------------------------------------------------------------------------------------

\end{document}  
